\subsection{环上的赋值}

设 \(\Gamma\) 是以加法记写的全序交换群。在本章的其余部分，对这样一个群，我们须考虑把一个元素添加到 \(\Gamma\) 上所得到的集，该元素记作 \(+\infty\)；把这个集记作 \(\Gamma_{\infty}\)，并赋予它：(1) 一个使 \(+\infty\) 为最大元素的全序，换言之，使得 \(\alpha < +\infty\)（对所有 \(\alpha \in \Gamma\)）；(2) 一个交换幺半群结构，其法则在 \(\Gamma\) 上诱导出给定的群法则，且由等式

\[
(+ \infty) + (+ \infty) = + \infty , \quad \alpha + (+ \infty) = + \infty
\]

（对所有 \(\alpha \in \Gamma\)）定义；立即可验证这个法则是结合的、交换的，且关系 \(\alpha \leqslant \beta\)（在 \(\Gamma_{\infty}\) 中）蕴涵 \(\alpha + \gamma \leqslant \beta + \gamma\)（对所有 \(\gamma \in \Gamma_{\infty}\)）。

定义 1. 设 C 是一个（不必交换的）环，\(\Gamma\) 是以加法记写的全序交换群。C 的以 \(\Gamma\) 为值群的任一赋值是指满足下列条件的任一映射 \(v: \mathbf{C} \to \Gamma_{\infty}\)：

\[
(\mathrm{VL}_{\mathrm{I}}) \quad v (x y) = v (x) + v (y) for x \in \mathbf {C}, y \in \mathbf {C}.
\]

\[
(\mathrm{VL}_{\mathrm{II}}) \quad v (x + y) \geqslant \inf (v (x), v (y)) for x \in \mathbf {C}, y \in \mathbf {C}.
\]

\[
(\mathrm{VL} _ {\mathrm{III}}) \quad v (1) = 0 and v (0) = + \infty .
\]

若 C 除 0 外没有零因子，则唯一的映射 \(v_0\)（从 C 到 \(\Gamma_\infty\)）——满足 \(v_0(x) = 0\)（对 \(x \neq 0\)）且 \(v_0(0) = +\infty\)——是一个赋值，称为 C 上的非真赋值。若 \(z \in \mathbf{C}\) 满足 \(z^n = 1\)（对某个整数 \(n \geqslant 1\)），则由 \((\mathrm{VL}_{\mathrm{I}})\)，\(nv(z) = v(z^n) = 0\)，从而 \(v(z) = 0\)（对 C 上的每个赋值 \(v\)），因为 \(\Gamma\) 是全序群。特别地 \(v(-1) = 0\)，于是 \(v(-x) = v(x)\)（对所有 \(x \in \mathbf{C}\)）。此外，由 \((\mathrm{VL}_{\mathrm{I}})\) 得出 \(v(xy) = v(yx)\)（对 C 中所有 \(x, y\)）。若 \(x\) 在 C 中可逆，则 \(v(x^{-1}) = -v(x)\)。

命题 1. 设 \(v\) 是 \(a\)（不必交换的）环 \(\mathbf{C}\) 上的赋值。对任意元素 \(x_{i} \in \mathbf{C}\)（\(1 \leqslant i \leqslant n\)），

\[
v \left(\sum_ {i = 1} ^ {n} x _ {i}\right) \geqslant \inf _ {1 \leqslant i \leqslant n} v (x _ {i})\tag{1}
\]

此外，若存在唯一的指标 \(k\) 使 \(v(x_k) = \inf_{1 \leqslant i \leqslant n} v(x_i)\)，则 (1) 的两端相等。特别地，若 \(v(x) \neq v(y)\)，则 \(v(x + y) = \inf(v(x), v(y))\)。

关系 (1) 由公理 \((\mathrm{VL}_{\mathrm{II}})\) 对 \(n\) 用归纳法得出。若存在唯一的指标 \(k\) 使 \(v(x_k) = \inf_{1 \leqslant i \leqslant n} v(x_i)\)，则记 \(y = \sum_{i \neq k} x_i\)、\(z = \sum_{i=1}^{n} x_i\)，由 (1) 有 \(v(y) > v(x_k)\) 且 \(v(z) \geqslant v(x_k)\)；若 \(v(z) > v(x_k)\)，则关系 \(x_k = z - y\) 将给出 \(v(x_k) \geqslant \inf(v(z), v(y)) > v(x_k)\)，这是荒谬的；于是 \(v(z) = v(x_k)\)，这就证明了第二个断言。

推论. 若 C 的元素的一个有限序列 \((x_{i})_{1\leqslant i\leqslant n}\)（其中 \(n\geqslant 2\)）满足 \(\sum_{i = 1}^{n}x_{i} = 0\)，则至少存在两个不同的指标 \(j,k\)，使

\[
v (x _ {j}) = v (x _ {k}) = \inf _ {1 \leqslant i \leqslant n} v (x _ {i}).
\]

若只有唯一的指标 \(k\) 使 \(v(x_k) = \inf_{1 \leqslant i \leqslant n} v(x_i)\)，则命题 1 将表明 \(v(x_k) = v(0) = +\infty\)，从而 \(v(x_i) = +\infty\)（对所有 \(i\)），这与关系 \(n \geqslant 2\) 及对 \(k\) 所作的假设相反。

注

\begin{enumerate}
\def\labelenumi{(\arabic{enumi})}
\item
  若 \(v: \mathbf{C} \to \Gamma_{\infty}\) 是 \(\mathbf{C}\) 上的赋值，\(u: \mathbf{B} \to \mathbf{C}\) 是从环 \(\mathbf{B}\) 到 \(\mathbf{C}\) 的同态，则立即可见合成映射 \(\mathbf{B} \xrightarrow{u} \mathbf{C} \xrightarrow{v} \Gamma_{\infty}\) 是 \(\mathbf{B}\) 的以 \(\Gamma\) 为值群的赋值。
\item
  条件 \((\mathrm{VL}_{\mathrm{I}})\) 与 \((\mathrm{VL}_{\mathrm{II}})\) 立即表明，集 \(\overline{v}^{-1}(+\infty)\) 是 C 中的一个双边理想 p，且由 \((\mathrm{VL}_{\mathrm{III}})\)，它异于 C；此外，若 x、y 是 C 的两个满足 \(v(xy)=+\infty\) 的元素，则由 \((\mathrm{VL}_{\mathrm{I}})\) 必有 \(v(x)=+\infty\) 或 \(v(y)=+\infty\)；换言之，商环 C/p 除 0 外没有零因子；立即可验证，由 v 通过取商导出的映射 \(\overline{v}:C/p\to\Gamma_{\infty}\) 是 C/p 上的赋值，且这个赋值下 \(+\infty\) 的原像约化为 0。
\end{enumerate}

